\documentclass[11pt]{article}

\usepackage[T1]{fontenc}
\usepackage{lmodern}
\usepackage{amsmath}
\usepackage{amssymb}
\usepackage[margin=1in]{geometry}
\usepackage{booktabs}
\usepackage{hyperref}

\newcommand{\key}[1]{\texttt{#1}}
\newcommand{\dd}{\mathrm{d}}

\title{Discrete diffusion Monte Carlo for Compton scattering:\\
the lattice, the interface, the energy update and the moments,\\
as they are to be implemented (design note, phase D0)}
\author{Athena++ Monte Carlo radiation transport}
\date{Draft of 2026-10-28, for approval before phase D1}

\begin{document}
\sloppy
\maketitle

\section{What the method does and where it lives}
\label{sec:what}

A photon inside an optically thick cell scatters many times before it leaves, and the
in-cell random walk of \texttt{mrw\_guide.tex} replaces those scatterings only when the
cell is many mean free paths thick, paying a dead zone at every face.  Discrete
diffusion Monte Carlo (DDMC; Densmore et al.\ 2007) replaces the photon in a thick cell
by a \emph{walker}: it has a cell, an energy, a weight and a time, and no
position or direction.  It waits in the cell for an exponential time whose rate is the
sum of the leakage rates through the six faces, then hops to the neighbour the chosen
face leads to; if that neighbour is thin, it emerges there as an ordinary transport
photon on the face.  A transport photon entering a thick cell is converted the other
way with the probability of the asymptotic diffusion limit.  The lattice reproduces the
diffusion equation on the grid to first order in the cell size, cell by cell, with each
cell's own properties, and costs one draw per cell traversed rather than one per
scattering.

For Compton scattering the Thomson extinction is grey, so the leakage rates do not
depend on the walker's energy and the energy decouples from the walk: after a dwell
of comoving path $ct$ the energy is drawn from the Green's function of the Kompaneets
equation at $y=\theta\,n_e\sigma_T\,ct$, exactly as the random-walk step does, with the
free-free survival from the same table (\texttt{kgreens.cpp};
\texttt{doc/monte\_carlo/kompaneets/}).  No energy groups exist.

The code is \texttt{src/monte\_carlo/ddmc.cpp} with \texttt{ddmc.hpp}, off unless
\key{<ddmc> enable} is set, and reaches the rest of the module through the hooks of
the plan's Section~2.9 and no others: a branch at the top of each pusher's per-photon
loop that runs the walk while the walker's flag is set, a conversion test where a
transport photon's \texttt{UpdateZone} enters a new cell, the integer photon column
\texttt{ddmcp} registered only when enabled, one guard in
\texttt{TransferPhotonsOnBlock}, one call in \texttt{SetupBlockFromFluid} that builds
the per-cell rate table, and the keys and report lines in \texttt{montecarlo.cpp}.
With the feature off, lists, moments and payloads are byte-identical to a build before
it existed.

\section{Conventions and units}
\label{sec:units}

The extinction $\chi$ is the comoving Thomson extinction in cm$^{-1}$, $\chi=n_e
\sigma_T$ with $n_e$ the electron density in \texttt{species(0,k,j,i)}, as the random
walk uses it; the free-free absorption enters only through the Green's function's sink
(Section~\ref{sec:energy}).  Lengths in the rate table are physical: cell widths in
code length times \key{l\_cgs}; under the general pusher the widths are the proper
lengths of the cell's coordinate edges in the cell-centre tetrad of the normal observer,
which \texttt{GeneralPusher::MRWFrame} builds (\texttt{fr.W}), and in the
spherical-polar pusher the box $\dd r\times r\,\dd\theta\times r\sin\theta\,\dd\phi$
that \texttt{CellGeometry} makes of the cell.  The diffusion coefficient is
$D=1/(3\chi)$, a length.  Times are comoving proper times inside the lattice; the dwell
is converted to coordinate time once per hop (Section~\ref{sec:gr}).  Rates are kept per
unit comoving path $c\,t$, in cm$^{-1}$, so that a leakage ``opacity'' $k_f$ competes
with the extinction on the same footing and the dwell is drawn as a path, $c\,t = -\ln
\xi/k_{\rm tot}$, with $\xi$ uniform.

\section{The lattice and its rates}
\label{sec:rates}

Let $E_c$ be the radiation energy density of cell $c$ in its comoving frame.  The
diffusion equation on the grid, in the finite-volume form of Densmore et al.\ (2007)
and Smith et al.\ (2018, their equation 10), is
\begin{equation}
  \frac{1}{c}\frac{\partial E_c}{\partial t}
  = \sum_{f} k_f^{(n\to c)}\,\frac{V_n}{V_c}\,E_n - \sum_f k_f^{(c\to n)}\,E_c
  - \kappa_{a,c}E_c ,
\end{equation}
where the sums run over the faces $f$ of the cell, $n$ is the neighbour across $f$, and
the leakage rate per unit path from $c$ into $n$ across a face normal to axis $i$ is
\begin{equation}
  k_f^{(c\to n)} = \frac{2D_c/\Delta_{c,i}}{\Delta_{c,i} + \Delta_{n,i}\,D_c/D_n}
  = \frac{2}{3\,\Delta_{c,i}\,(\chi_c\Delta_{c,i} + \chi_n\Delta_{n,i})} .
  \label{eq:kint}
\end{equation}
This is Fick's law through the two half-cells, $F = -(E_n - E_c)/\bigl(\Delta_c/2D_c +
\Delta_n/2D_n\bigr)$, with the part proportional to $E_c$ divided by the cell width.
In a uniform medium it is $1/(3\chi\Delta^2)$, and a walker makes about
$(\chi\Delta)^2/3$ Thomson paths' worth of dwell per hop, which is the number of
scatterings a hop replaces.

A face whose far side is not a diffusion cell (a transport cell, a face to another
refinement level under the rule of Section~\ref{sec:faces}, or a mesh boundary) is a
\emph{boundary face}.  Its rate comes from the asymptotic diffusion-limit boundary
condition of Habetler \& Matkowsky (1975), which Densmore et al.\ (2007) and
Abdikamalov et al.\ (2012, their equations 120 to 127) use in the same way: at the face
the mean intensity obeys
\begin{equation}
  \int_0^1 W(\mu)\,I_b(\mu)\,\dd\mu = J_{\rm face} - \frac{\lambda}{\chi}
  \left.\frac{\partial J}{\partial n}\right|_{\rm face},\qquad
  W(\mu)\simeq \mu + \tfrac32\mu^2,\qquad \lambda = 0.7104,
  \label{eq:abc}
\end{equation}
with $I_b$ the intensity incident from outside, $n$ the outward normal and
$\int_0^1 W\,\dd\mu = 1$.  Replacing the derivative by $(J_{\rm face} - J_c)/(\Delta_c/2)$
and the face flux by Fick's law over the half-cell, $H = -(2/3\chi_c\Delta_c)(J_{\rm face}
- J_c)$, gives
\begin{equation}
  H = \frac{2}{3\chi_c\Delta_c + 6\lambda}\Bigl(J_c - \int_0^1 W(\mu)\,I_b(\mu)\,\dd\mu\Bigr),
\end{equation}
so the leakage rate out through a boundary face is
\begin{equation}
  k_f^{\rm (bdy)} = \frac{2}{3\,\Delta_{c,i}\,(\chi_c\Delta_{c,i} + 2\lambda)}
  = \frac{2D_c/\Delta_{c,i}}{\Delta_{c,i} + 2\lambda/\chi_c} ,
  \label{eq:kbdy}
\end{equation}
the interior form with the far half-cell replaced by the extrapolation distance
$\lambda/\chi_c = 2.13\,D_c$, and the incident intensity enters as a source with
weight $W(\mu)$, which is the conversion probability of Section~\ref{sec:interface}.
(The plan of 2026-10-06 had $2\lambda D_c$ in the denominator, a factor three short.)

A cell is a diffusion cell when
\begin{equation}
  \tau_{\rm cell} = \chi_c \min_i \Delta_{c,i} \ge \tau_{\rm ddmc}
\end{equation}
with \key{<ddmc> tau} the threshold; the thinnest dimension governs because leakage
is fastest across it.  Densmore et al.\ validate the interface down to $\tau_{\rm cell}$
of about 2 and Smith et al.\ suggest 5; D1 scans 2, 3, 5 and 10 against the analog run
and sets the default.  One threshold serves both directions; the asymmetry between
entering and leaving is carried by the probabilities, and a second key is added only if
D1 asks for it.

The rate table is built once per transport cycle in \texttt{SetupBlockFromFluid},
after the number densities, as \texttt{ddmc\_rate(7,k,j,i)}: six face rates and
their sum, over the active cells, with the neighbour's $\chi_n$ and $\Delta_n$ read from
the ghost zones.  Ghost zones are valid for this: \texttt{Mesh::Initialize}
(\texttt{mesh.cpp}, the block from the conserved-variable exchange through
\texttt{ProlongateBoundaries} to \texttt{ConservedToPrimitive} over the extended
bounds and \texttt{ApplyPhysicalBoundaries}) fills the primitives in the ghost zones
from the neighbours, across levels, before \texttt{MonteCarlo::Initialize} runs
(\texttt{main.cpp}), and the former skip of the inversion for GR Monte Carlo builds is
gone; \texttt{FillBounds} then spans the ghosts when \texttt{rho}, \texttt{tgas} and
\texttt{species} are filled.  A problem generator's \texttt{ExcludeCells} must zero
the species over the same bounds, which D1 asserts on a two-block run by comparing a
ghost $\chi$ with the neighbour's active value.  Under the general pusher the widths
need a tetrad per cell, including the ghost layer; the per-cell transformations of
\texttt{ComputeTransformations} are where that goes.  The table is rebuilt on restart
and redistribution like the opacity tables, so nothing is carried.

\section{The walker and its life}
\label{sec:walker}

A walker is a \texttt{Photon} whose integer column \texttt{ddmcp} is 1.
Its position is the centre of its cell, which keeps every consumer of a position
working unchanged: \texttt{UpdateZone}, the boundary functions, the block exchange, the
outputs, the load balancer's payload and the restart record.  Its direction columns
are unused and its coherency tensor and Stokes parameters are zero.  Its energy
\texttt{ep} is the comoving energy in its cell's frame; its weight, time
\texttt{x0p}, budget \texttt{dtp}, scattering count \texttt{nscp} and type are what
they are for a photon.  Its flight columns (\texttt{taup}, \texttt{strp}) are reset
at emergence.

The walk, in the pusher's per-photon loop while \texttt{ddmcp} is set:
\begin{enumerate}
\item Read the cell's rates.  Form the path budget $c\,t_b$: the cycle budget
  \texttt{dtp} converted to comoving path (Section~\ref{sec:gr}), and the Kompaneets
  table's $y$ range, $c\,t \le y_{\max}/(\theta n_e\sigma_T)$.  (The random walk's
  \key{accel\_pmax} has no analogue; the lattice has no domain survival.)
\item Draw the path $c\,t = -\ln\xi/k_{\rm tot}$.  If it exceeds $c\,t_b$, set $c\,t =
  c\,t_b$ and stay in the cell (the walker keeps its flag; a budget stop returns it
  to the block loop, a table cap continues the walk).
\item Update the energy and weight over the dwell (Section~\ref{sec:energy}); deposit
  the moments and source terms (Section~\ref{sec:moments}); advance \texttt{x0p} and
  \texttt{dtp}; add $n_e\sigma_T\,ct$ to \texttt{nscp} and to the report's count of
  scatterings replaced.
\item If a budget stop or absorption ended the dwell, return.  Otherwise choose the
  face with probability $k_f/k_{\rm tot}$ and hop:
  \begin{itemize}
  \item interior face: move the walker to the neighbour's centre, transform its
    energy into the neighbour's frame (Section~\ref{sec:moving}), deposit the face
    counters (Section~\ref{sec:moments}); if the neighbour is a ghost cell, the walker
    is now outside the block and \texttt{UpdateZone} buffers it to the neighbour
    block, where it lands in the face cell with its flag and continues;
  \item boundary face: emergence (Section~\ref{sec:interface}), which clears the flag;
    the loop then continues as a transport photon with a fresh optical depth, and a
    mesh boundary is met by the ordinary boundary function on the first step.
  \end{itemize}
\end{enumerate}
The general pusher invalidates its cached metric and connection after a hop, as the
random-walk branch does.  \texttt{UserWorkInMove} is called after each hop, so a
generator's hook (the horizon test of \texttt{mc\_readhdf\_gr}) sees the walker at
its new centre.  \texttt{TransferPhotonsOnBlock} skips absorption and scattering for a
walker that returns from \texttt{Move} with the flag set.

\section{Faces that are not interior}
\label{sec:faces}

\textbf{Mesh boundaries.}  A face on the mesh boundary is a boundary face whatever the
boundary condition: the walker emerges through it and the boundary function
(escape, absorb, periodic, polar) acts on the transport photon as it does today.  A
periodic neighbour is in fact a real cell whose ghost values are filled, so periodic
faces could be interior; they are left as boundary faces in the first build for
simplicity and the report counts them.

\textbf{Refinement levels.}  A face whose neighbour block is at another level (known
from \texttt{pbval->nblevel} for the block's outer layer of cells; interior faces of a
block never jump levels) is a boundary face in both directions: a walker reaching it
emerges through it, a transport photon crossing it into a thick cell converts by the
usual probability.  This is exact and loses efficiency only at level jumps; the report
counts hops stopped there, and the split-face rates of Densmore et al.\ (2007), which
weight the four fine sub-faces by area, are the follow-up if the count is large on the
XRB, whose snapshot has eight levels.

\textbf{Block faces at the same level} are interior faces: the rate uses the ghost
neighbour, and the hop is the BUFFERED exchange.

\section{Energy and absorption over a dwell}
\label{sec:energy}

With $\theta = kT_c/m_ec^2$ of the cell and the walker's comoving energy $x_i = E/kT_c$,
a dwell of path $ct$ has Compton parameter $y = \theta\,n_e\sigma_T\,ct$ and the new
energy is $x_f kT_c$ with $x_f$ from \texttt{kgreens->Quantile(xi, y, lam, u)}; the
weight is multiplied by \texttt{kgreens->Survival(xi, y, lam)} when the absorption is
free-free (\texttt{lam} the cell's free-free parameter as \texttt{MRWStep} forms it from
\texttt{ff\_cell}), and by $\exp(-\kappa_a\,ct)$ for a grey absorption that is not
free-free.  Under \key{abs\_method = prob} the survival is a probability instead and the
walker is ABSORBED with its complement; under \key{weight} it is a weight and the
roulette of \texttt{RouletteOrAbsorb} applies below \key{minweight}.  The guards are
those of \texttt{MRWTrigger} and \texttt{MRWStep}: $\theta\le0.03$, $x\theta\le0.05$,
$x$ inside the table, the sink strength within the table's range, $n_e>0$; a walker
whose cell fails them is not converted (a transport photon in such a cell stays a
transport photon, and the conversion test includes the guards), and the report counts
the declines by reason.  A long dwell is capped at the table's $y_{\max}$ and the walk
continues with a second draw, so no energy is ever extrapolated.

Only the Thomson-grey regime is covered; a Klein-Nishina dependence of the leakage
rates on energy would be a per-group rate table, which the XRB bands below 10 keV do
not need.

\section{Moments and source terms}
\label{sec:moments}

Required: energy densities, fluxes and the radiation force must come out of diffusion
cells.

\textbf{Energy density}, per dwell, by the time-spent estimator:
\texttt{UpdateMomentsMRW(pphot, ct, $\bar E$, beta\_tet, gam, ip)} with $\bar E$ the
mean of the energies before and after the dwell, which deposits the isotropic comoving
field $E_c = w\bar E\,ct/c$ boosted to the lab (its Section~4 formulas).

\textbf{Flux}, by the face-hop estimator.  Each hop through a face normal to axis $i$
adds $s\,wE$ to that face's counter of each of the two cells sharing it, $s=\pm1$ the
direction of crossing along $+i$ and $E$ the energy carried across (the leaving cell's
comoving energy; the frame change of the hop is second order in the flux).  Over an
integration time the counter divided by the face area and \texttt{tint} is the flux
through the face, and the cell-centred comoving flux along $i$ is the mean of the
cell's two faces.  It is the discrete Fick flux of the lattice, $F=-cD\nabla E$ on the
grid, so it is consistent with what the lattice solves and needs no model.  The
counters are a per-block array \texttt{ddmc\_flux(6,k,j,i)} in the comoving frame;
both sides of a block face record every crossing themselves (the leaving block when the
walker leaves, the receiving block when it arrives), so no communication is needed.
At normalization the cell's comoving $(E_c, F_c, P_c = E_c\delta/3)$ is boosted into
the lab moments with the cell's velocity matrices of \texttt{ComovingFrameMatrix}, the
isotropic part having been deposited per dwell already, so only the flux term is
added there.

\textbf{Pressure}: $E_c/3$ on the diagonal in the comoving frame, the diffusion
closure, which is the assumption of a diffusion cell and not a measurement.

\textbf{Radiation force.}  The momentum deposited in the gas of a diffusion cell per
unit volume and time is $\chi_c F_c/c$ with $\chi_c$ the total extinction, equal to
$-\nabla E_c/3$ by Fick's law.  With
the flux from the face counters, each hop deposits
\begin{equation}
  \delta p_i = s\,\frac{wE\,\chi_c\,\Delta_{c,i}}{2c}
\end{equation}
into \texttt{sourceterms(MCRS}$_i$\texttt{)} of each cell sharing the face, each with
its own $\chi$ and width (the factor $1/2$ is the face average), and the same into
\texttt{MCRF}$_i$, since the flux force and the momentum force coincide on the
lattice.  The sign convention is that of \texttt{UpdateSourceTerms}, where the gas
gains minus the photon's momentum change.  \textbf{Energy exchange}
\texttt{MCRS0}: $w_{\rm after}E_{\rm after} - w_{\rm before}E_{\rm before}$ per dwell,
through \texttt{UpdateSourceTerms} with the energy and weight before and after, which
also covers the absorbed energy.  The per-walker contributions that
\texttt{sourceterms\_error} holds back until the walker leaves a cell accumulate over
its hops in that cell unchanged.

\textbf{Known limits.}  The flux in a thick cell is the difference of two large hop
counts, with a statistical error set by the gross hop rate; the analog flux there is the
residual of near-isotropic flights and is no better.  The scattering-binned spectral
moments (\texttt{Jnu}, \texttt{mom\_flag\_scat}) see only the energies before and
after a dwell.  Fick's law from the energy-density estimator of neighbouring cells is a
cross-check in the tests, not the product, since it fails in interface cells where
transport flights also contribute.

\section{The interface}
\label{sec:interface}

\textbf{Conversion}, when a transport photon reaches a face of its cell beyond which
lies a diffusion cell $n$, with direction cosine $\mu$ to the face normal.  From
equation~(\ref{eq:abc}) the incident intensity enters the diffusion solution with weight
$W(\mu)$, and the part of the face flux it drives is $4\pi\,(2/(3\chi_n\Delta_{n,i} +
6\lambda))\,\langle W I_b\rangle$, where $\langle W I_b\rangle = \int_0^1 W I_b\,\dd\mu$
is an average normalized like the mean intensity, over the full sphere, while the
energy a beam carries through the face is $\mu I_b\,\dd\Omega$ with the hemisphere's
$2\pi$.  The ratio is the conversion probability,
\begin{equation}
  P(\mu) = \frac{4}{3\chi_n\Delta_{n,i} + 6\lambda}\Bigl(1 + \tfrac32\mu\Bigr)
  = \frac{4}{3\,(\chi_n\Delta_{n,i} + 2\lambda)}\Bigl(1 + \tfrac32\mu\Bigr),
  \label{eq:pconv}
\end{equation}
Smith et al.'s equation 18, $2(2/3+\mu)/(\chi\Delta + 2\lambda)$, and the form of
Densmore et al.\ (2007); it derives from the same finite difference as the boundary
leakage rate~(\ref{eq:kbdy}), and the two must be used together.  The check that
fixes the factor: a diffusion cell in equilibrium with an isotropic field $I_0$
outside leaks $4\pi I_0\cdot 2/(3(\chi\Delta+2\lambda))$ per unit area and receives
$2\pi I_0\int_0^1 \mu P(\mu)\,\dd\mu = 8\pi I_0/(3(\chi\Delta + 2\lambda))$, equal
only with the 4.  (The draft of this note had a 2, from counting the hemisphere
twice; built that way, the lattice lost energy at every interface and the Thomson
sphere escaped 5 to 13 percent too fast, while a cube with no interface reproduced
its exact lattice solution.  The factor was found on 2026-10-28 by that pair of
tests, the sphere's error vanishing once it was corrected.)  $P$ exceeds 1 below
$\chi\Delta \simeq 1.9$ at normal incidence, which is why the threshold must stay at
2 or above; a value above 1 is clamped and counted.  A converted photon becomes a
walker at the cell's centre with its energy transformed into the cell's frame.  A
photon not converted is returned into the cell it came from with an isotropic inward
direction, at the face, keeping its flight state.

\textbf{Emergence}, when a walker hops through a boundary face.  The photon is placed
at a uniform point on the face, on the transport side, with the flag cleared, an energy
drawn in the new frame (Section~\ref{sec:moving}), weight carried, Stokes parameters
zero, \texttt{taup} $=-1$ and \texttt{strp} $=0$ so the pusher draws a fresh flight.
Its outward direction is drawn from one of two distributions, \key{<ddmc> emergence}:
\key{asymptotic}, $W(\mu) = \mu + \tfrac32\mu^2$ on the outward hemisphere, the
flux-weighted emergent intensity $I\propto\tfrac23+\mu$ of the diffusion limit and the
distribution consistent with equation~(\ref{eq:abc}); or \key{isotropic}, uniform in
$\mu\in(0,1)$, which is what Densmore et al.\ (2007), Abdikamalov et al.\ (2012) and
Smith et al.\ (2018) use.  The first is the expected default; the disk atmosphere
against Feautrier decides in D1, since the emergent intensity versus $\mu$ is the one
observable that sees this choice.  $W(\mu)$ is sampled by solving
$\mu^2/2 + \mu^3/2 = \xi$ for $\mu$, a few Newton steps from $\mu=\xi^{1/2}$.

\section{A moving medium}
\label{sec:moving}

The lattice is in the comoving frame of each cell: rates, dwell, energy and moments.
At a hop the energy changes frame.  Rather than a first-order Doppler formula with a
mean cosine (Smith et al., their equations 19 and 20), the hop samples a direction from
the emergence distribution about the face normal, builds the comoving wavevector of
energy $E$ in the leaving cell, transforms it to the lab with that cell's velocity
(\texttt{TransformToCoordinate} at the cell centre), moves to the entering cell's
centre, transforms it into that cell's frame (\texttt{TransformToComoving}), keeps the
energy and discards the direction.  This is exact in distribution for the diffusion
limit, uses the transforms that exist, and costs a few operations per hop.  The dwell's
lab time is $\gamma$ times the comoving time (Section~\ref{sec:gr}).

What is not represented: the advection of the fluid past the lattice within a cell, and
velocity gradients within a cell.  The lattice is the grid and the cell's velocity is
one number; this is the approximation Smith et al.\ make and the moving sphere of D4
measures (beta 0.01, 0.1, 0.3 against analog).  The random walk's advection room has no
analogue.

\section{Curved spacetime}
\label{sec:gr}

Under the general pusher the cell's widths are the proper lengths of its coordinate
edges in the normal observer's tetrad at the cell centre (\texttt{fr.W} of
\texttt{GeneralPusher::MRWFrame}), $\theta$ and $n_e$ are the cell's comoving values,
and the dwell is a comoving proper time $t$.  Coordinate time advances by
$\Delta x^0 = \gamma\,ct/(c\,\alpha)$ in code units with $\gamma = u^{\hat t}$ the
fluid's Lorentz factor on the tetrad and $\alpha$ the lapse, exactly as
\texttt{MRWStep} does (\texttt{dl = gam*ct/l\_cgs/fr.lapse}); \texttt{dtp} decreases by
the same.  The hop transform of Section~\ref{sec:moving} goes through the two cells'
tetrads at their centres, which covers the gravitational shift as well as the Doppler
one.  The walker's coordinate position is the cell centre.  Inside the horizon of
\texttt{mc\_readhdf\_gr} the species are zero and no cell is a diffusion cell.

\section{Time and dynamic runs}
\label{sec:time}

The waiting times are physical.  In a static run \texttt{tmax} is huge and no budget
binds.  In a dynamic run the budget is \texttt{dtp}, and a walker whose dwell would
exceed it stops in its cell with the flag set and returns to the block loop.

Finding of D0, by reading the three pushers, \texttt{TransferPhotonsOnBlock} and
\texttt{FinishRound}: today \texttt{dtp} is only the exit condition of the pusher's
step loop.  A transport photon that exhausts it returns EVOLVING, is scattered in place
by the block loop at every pass, and keeps the round from finishing, since the round
ends only when no photon is on any list.  So the budget is not a stop today, and
dynamic runs work because \texttt{tmax} (the hydro step) does not bind on their
photons.  A true stop, for photons and walkers alike, needs a held status that the
round accounting treats as done and the next cycle resumes, which is outside the
plan's hook list and is a decision for the user before D6.  Until then the walker's
budget stop is implemented as the photon's is (the walker sits, flagged, and is
skipped by the block loop), and the Compton phases run static.

\section{Keys and report lines (proposed)}
\label{sec:keys}

\begin{center}
\begin{tabular}{llp{8.2cm}}
\toprule
key (\texttt{<ddmc>}) & default & meaning \\
\midrule
\key{enable} & false & turn the lattice on \\
\key{tau} & 5 (D1 sets it) & cell optical depth, extinction times the smallest width, above which a cell is a diffusion cell \\
\key{emergence} & \texttt{asymptotic} & direction distribution at emergence: \texttt{asymptotic} ($W(\mu)$) or \texttt{isotropic} \\
\key{report} & false & print the counters below per transport \\
\key{kgreens\_file} & (none) & the Kompaneets table under Compton scattering, read from here or from \texttt{<randomwalk>} (D2) \\
\bottomrule
\end{tabular}
\end{center}

Nothing else is read from \texttt{<montecarlo>} or \texttt{<randomwalk>}; the random walk
never fires in a run with the lattice on (its threshold lies above \key{tau} and a
walker has no position), and the wiki says the two do not combine.  The report
prints: photons converted, photons returned, walkers emerged, hops, hops through
block faces, hops stopped at level jumps and mesh boundaries, scatterings replaced,
budget stops, table caps, declines by reason (the Kompaneets guards), and the
\key{accel\_report} histogram of analog scatterings by cell thickness.

\section{What D1 built, where it departs from the sections above}
\label{sec:d1}

Phase D1 (2026-10-28, branch \texttt{ddmc}): grey elastic scattering in flat space in
the Cartesian and spherical-polar pushers, one file \texttt{ddmc.cpp} plus the hooks of
Section~\ref{sec:what}.  The departures from the design, each for a reason found in
building or testing:
\begin{itemize}
\item \textbf{Conversion at the start of a flight.}  A transport photon beginning a
  fresh flight after a scattering in a diffusion cell is converted with probability 1
  at the cell's centre.  A photon emitted in a diffusion cell flies its first flight
  before it can be converted, so a direct escape from the top layer of the lattice
  keeps the position it was born at; the walker's centre placement is the one
  $O(\Delta)$ error the method has for emergent light from an emitting medium, and it
  falls with the threshold as the light from the first lattice layer does.  Emission
  inside diffusion cells therefore makes walkers already, which the plan had in D6;
  a photon that reached a diffusion cell through an edge or corner crossing, or across
  a block face from a neighbour that could not test it, is caught the same way at its
  next scattering.
\item \textbf{The face test precedes the crossing.}  \texttt{DDMCFaceTest} runs on the
  target cell before \texttt{MovePhotonToNextZone}, so a ghost cell across a same-level
  block face is tested too and a conversion into it is buffered to that block.  A
  returned photon and an emerged photon are placed $10^{-7}$ of the cell's width inside
  the cell (the reflecting boundaries' device), so the crossing is the pusher's own
  first step and every boundary function acts as usual.
\item \textbf{The flag carries the arrival.}  \texttt{ddmcp} is 1 for a resident
  walker and $2 + f$ for one that has just hopped (or converted) through a block face
  and enters its new cell through face $f$: the receiving block then makes the
  entering deposit and the frame change, so both sides of a block face record every
  crossing without communication.  The lab energy and direction ride along in the
  unused columns while the walker is in transit.
\item \textbf{The mask over the ghosts.}  \texttt{ddmc\_cell} covers the ghost zone,
  real only across a same-level block face (\texttt{block\_bcs} and \texttt{nblevel});
  the problem generator's \texttt{accel\_mask} is shared with the random walk and
  \texttt{mc\_sphere\_isoth} now fills it over the ghosts, so two blocks agree on a
  cell.  A walker whose cell is no longer a diffusion cell (a rebuilt lattice, or a
  block that disagrees with its neighbour's ghost) is released as an isotropic transport
  photon at the centre, counted.
\item \textbf{Moments.}  The hop flux goes straight into the lab moments (and the
  comoving ones when they are accumulated), with a first-order boost when the cell
  moves, rather than through a separate face array; the energy density of a dwell
  uses the mean weight along it under absorption, $w(1-e^{-\kappa_a ct})/(\kappa_a
  ct)$; the scatterings a dwell stands for are counted up to the absorption in
  expectation.  The pending force-error contribution is flushed at each change of
  cell, so a walker's visits to a cell count separately.
\item \textbf{Rates with absorption.}  The table holds the pure Thomson rates; a
  walker with an absorption opacity forms its rates, its conversion probability and
  its force from the total extinction at its energy (Section~\ref{sec:rates}).
\item \textbf{Threshold.}  \key{tau} defaults to 3 and is refused below 2, where
  $P(1)$ exceeds 1 (Section~\ref{sec:interface}); a clamped probability is counted.
\item \textbf{Report.}  Cells, dwells, scatterings stood for, hops and hops through
  block faces, conversions at faces and at flight start, returns, emergences, releases,
  budget stops, clamps; \key{report} also turns on the histogram of analog
  scatterings by cell depth.
\end{itemize}

\textbf{What the tests found.}  The conversion probability was the one error of the
draft (Section~\ref{sec:interface}), and only the pair of tests separated it from the
lattice's discretization: the Thomson sphere, whose interface shell of cut cells is
transport, escaped 5 to 13 percent too fast at every resolution, while the cube that
fills the mesh, where the lattice has no interface, reproduced its exact first-passage
solution (mean path and dwell count) to the statistics.  With the factor corrected the
sphere agrees with the analog run in the mean escape path to 0.5 percent at 16 and 32
cells across for cell depths from 2 to 16, and to 2 to 3 percent at 8 cells across
whatever the depth (the blocky sphere and its shell), the escape-path distributions
within the two-sample Kolmogorov scale.  Free-free absorption, monochromatic and so
grey, agrees under the weight method to $3\times10^{-4}$ in the escaping weight at an
absorbed fraction of 18 percent and under the probability method within the binomial
error; a case where 99.8 percent is absorbed, an exponential of the path, converges
from 30 to 6 percent between 16 and 32 cells.  In the cube's moments (depth 5), the
energy density agrees per cell to 5 percent with the lattice's expected shape (4
percent high at the source, 3 to 6 low at the boundary), the flux and the force to a
percent on average with the luminosity through the inner shells conserved, and
$\chi F/c$ equals the deposited force to half a percent in both runs; the off-diagonal
pressure is zero on the lattice.  A dwell costs 0.24 to 1 microseconds against 0.4 to
0.5 for a scattering, so the gain is the scatterings a dwell stands for: 100 times at
depth 10, 7 to 19 times on the gridded spheres at depths 5 to 8.

The disk atmosphere of \texttt{tst/montecarlo/disk\_atmosphere} (spherical polar, the
legacy pusher, polarized, free-free emission and absorption, cells up to 760 deep at
the midplane; the lattice takes 26 of its 131 thousand cells at a threshold of 3):
the ladder of 1, 4 and 16 million photons passes the Feautrier gate with the lattice
as without it (grazing $Q/I = -0.0395 \pm 0.0090$ against $-0.0401$, the norm falling as
$N^{-0.49}$), the emergent intensities agree bin by bin within their errors, and in
the moments (4 intervals of 500 thousand photons, the error from the scatter of the
interval sums, which carries the correlations along a photon's path) the energy
density and the vertical flux summed over the thin layers agree to 1 to 2 percent
within errors of 2.5 percent, the deep energy density to about a percent.  The deep
fluxes and forces are noise in both runs: where absorption dominates, a photon dies
after a scattering or two while a hop carries $\chi\Delta/2$ times a scattering's
momentum, so the lattice's force estimator is the noisier there by orders of
magnitude, a property to keep in mind for radiation forces in absorption-dominated
thick cells.  It was this test that showed the rates needed the total extinction.

\section{What D2 built, where it departs from Section~\ref{sec:energy}}
\label{sec:d2}

Built 2026-10-29 on branch \texttt{ddmc}, in \texttt{ddmc.cpp} with one check lifted in
the constructor.  The dwell under \key{scattering = compton} draws the walker's energy
from the Kompaneets table at $y = \theta\,N$, with $N$ the scatterings the dwell stands
for (the walker's scattering opacity times $ct$), in chunks of the table's $y_{\max}$
with a fresh draw per chunk (the process is Markov in $y$ at fixed $\theta$), and its
free-free survival from the same table with \texttt{lam} rescaled to the walker's $y$;
any other absorption multiplies by $\exp(-\kappa_a ct)$.  The opacities are refreshed at
the new energy, the source terms take the energy change as they take the absorbed
weight, and the moments use the mean of the energies before and after.  The table file
is \key{kgreens\_file} in \texttt{<ddmc>} or \texttt{<randomwalk>}.  Two counters were
added: conversions declined because the photon's energy fails the guards, and walkers
released because a dwell took their energy outside the table; the report prints both
under Compton.

Two departures from Section~\ref{sec:energy}, both forced by the Compton sphere.
First, there is no bound on $x\theta$.  With the random walk's $x\theta\le0.05$, at
$\theta = 0.01$ the lattice came out 6 percent low in the mean escaping energy with the
spectrum empty above $x = 5$: an energy at which walkers are released and transport
photons are not converted is a boundary in energy between dwells that stand for about
twenty scatterings and single scatterings, and a random walk with uneven steps across a
line depletes the side with the longer steps.  With the bound at $x\theta\le0.1$ the
mean energy agreed with the Kompaneets theory to 0.1 percent, so the lattice keeps only
the table's $x$ range, which for the cells that $\theta\le0.03$ admits lies far out in
the Wien tail, and the free-free sink guard.  Second, the leakage rates, the conversion
probability and $y$ take the walker's scattering opacity, the thermally averaged
Klein-Nishina value that the analog run uses, through the per-walker rates
(\texttt{DDMCRates}, \texttt{DDMCNeighbourExtinction}): with the Thomson rates the
lattice reproduced the Thomson diffusion theory's escape paths exactly while the analog
run's were 7 to 12 percent shorter at $\theta$ of 0.01 to 0.02, and its spectrum lacked
the hardening that the energy-dependent escape gives.  The cell guard $\theta\le0.03$ is
applied in \texttt{BuildDDMCRates}, so such cells are transport cells.

Measurements, the Compton sphere of \texttt{sphere\_compton.py} with $10^5$ photons, the
lattice against the driver's Green's-function gate and against the analog run of the
same deck ($x_0 = 1$, $\theta = 0.002$ unless said): the gate passes on the lattice as
on the analog run at $\tau$ of 15 (8 cells across, 57 diffusion cells), 40 (8 and 16
cells) and 100 (32 cells, 7249 diffusion cells), at $x_0$ of 0.1 and 3, and with
free-free absorption at a radius where it takes 14 percent (escaping weight 0.8606
against 0.8602 analog and 0.8579 theory at $x_0 = 0.3$, 0.9742 against 0.9741 at
$x_0 = 1$); the lattice and the analog spectra agree within 0.5 to 1.8 times the
two-sample Kolmogorov scale and the mean escaping energies to a percent, the escape
paths within 2 times the scale at 16 and 32 cells and 6 times at 8 cells (the coarse
lattice's short paths of Section~\ref{sec:d1}).  At $\theta$ of 0.01 and 0.02 with
both departures built, the lattice agrees with the analog run to 0.5 percent in the
mean energy and 1 to 1.4 times the scale in the spectrum, while both sit 2 to 3 percent
above the Thomson Kompaneets theory, which is the exact kernel's hardening and not the
lattice's.  At $\tau = 100$ on 32 cells the analog run's $4.8\times10^8$ scatterings
become $4.8\times10^6$ scatterings and $1.6\times10^7$ dwells.  The driver's own
free-free survival check fails identically on the analog and the lattice at the default
radius of $10^{10}$ cm, where the sink is $10^{-4}$: a deck limitation, which the radius
of $10^7$ cm used here avoids.

\section{What D3 built}
\label{sec:d3}

Built 2026-10-29.  The hop across block faces and the level-jump rule of
Section~\ref{sec:faces} were already in place from D1, so D3 added what the plan's
decision needs: \texttt{ddmc\_jump}, the bits of each active cell's faces that lead to a
block at another refinement level, the count of diffusion cells with such a face, and
the count of emergences through them, printed in the report as ``$N$ cells ($M$ at level
jumps)'' and ``emerged $E$ ($J$ through level jumps)''.  The mask and its counts are built
with \key{report = true} alone, without \key{enable}, so a run under the general pusher,
where the lattice is refused until D5, still reports what the lattice would reach.

Tests, in \texttt{tst/montecarlo/sphere\_compton/ddmc\_blocks.py}, on the Thomson
sphere with the lattice on: over 8 and 64 blocks on 1, 4 and 16 ranks every run
conserves its photons and agrees with the 8-block one-rank run within 1.1 times the
two-sample Kolmogorov scale in the escape paths (statistical, since each block seeds its
own generator); on one rank, plain against \texttt{<loadbalancing> test\_repack = local}
and \texttt{mesh} gives byte-identical lists, which is the check that a walker survives
the pack and unpack and that the lattice is rebuilt from the fluid on arrival; with a
static level-1 block of side $0.5R$ inside the sphere, the lattice agrees with the analog
run on the same mesh within 0.7 times the scale, with 318 of its 6511 diffusion cells
touching the jump and 121 thousand of its 229 thousand emergences going through it.
That last number is what Densmore's split-face rates would save on such a mesh, and the
decision to build them waits on the same count from the X-ray binary's mesh, for which
three report-only decks were written (threshold 2, 3 and 4).

The test found a bug in the generator, not in the exchange: \texttt{mc\_sphere\_isoth}
kept the indices of the cell holding the origin in file-scope variables that every
block's problem generator overwrote, so the emitting block used the last block's values.
With 8 blocks the last block holds the origin and nothing showed; with 27 or 64 the
point source was born with wrong indices, sent to the wrong neighbour and destroyed on
arrival, 20 to 50 percent of the photons under either pusher.  The origin cell is now
found on the emitting block in \texttt{InitializePhoton}.

\section{Decisions taken before D1}
\label{sec:decisions}

Settled with the user on 2026-10-28.
\begin{enumerate}
\item The emergence distribution: both \key{asymptotic} and \key{isotropic} are built
  behind the key, and the disk atmosphere against Feautrier in D1 picks the default
  (Section~\ref{sec:interface}).  No choice is made before the test.
\item Level jumps: a face whose neighbour block is at another refinement level is a
  boundary face in both directions (a walker emerges through it, a transport
  photon converts on the far side by the usual probability).  Exact, and only costs a
  conversion per crossing; the split-face rates of Densmore et al.\ (2007) are the
  follow-up if the report shows many hops stopped there on the XRB
  (Section~\ref{sec:faces}).
\item The dynamic budget: a held status that the round accounting treats as done and
  the next cycle resumes is wanted, for transport photons as much as for walkers,
  so it is a feature of its own outside the lattice's hook list and not a lattice
  phase.  Until it exists the walker's budget stop copies the photon's
  (Section~\ref{sec:time}), and the Compton phases run static.
\item The threshold is the extinction times the cell's smallest width, one key
  \key{tau} (Section~\ref{sec:rates}).  A cell thick in two directions and thin in the
  third is not a diffusion cell, since the diffusion limit fails along the thin one.
\item Periodic faces are boundary faces in the first build; the report counts the
  emergences through them (Section~\ref{sec:faces}).
\item The lattice runs only with \key{<ddmc> enable = true}, and the \texttt{ddmcp}
  column is registered on the \texttt{Photon} class only then, so a run with the
  feature off has the same photon layout, memory and output bytes as before the
  feature existed.
\end{enumerate}

\section{References}

Densmore, J. D., Urbatsch, T. J., Evans, T. M., Buksas, M. W. 2007, J. Comput. Phys.
222, 485.  Abdikamalov, E., Burrows, A., Ott, C. D., L\"offler, F., O'Connor, E.,
Dolence, J. C., Schnetter, E. 2012, ApJ 755, 111 (arXiv:1203.2915; the asymptotic
boundary condition, equations 120 to 128, and the conversion rule).  Smith, A., Tsang,
B. T.-H., Bromm, V., Milosavljevi\'c, M. 2018, MNRAS 479, 2065 (arXiv:1709.10187).
Habetler, G. J., Matkowsky, B. J. 1975, J. Math. Phys. 16, 846.  Wollaeger, R. T.,
van Rossum, D. R., et al. 2013, ApJS 209, 36.

\end{document}
